\documentclass[12pt,a4paper]{article}
\usepackage[utf8]{inputenc}
\usepackage[T1]{fontenc}
\usepackage[english]{babel}
\usepackage{geometry}
\geometry{margin=2.5cm}
\usepackage{amsmath,amssymb,graphicx,booktabs,array,microtype,xcolor}
\usepackage{xurl}
\usepackage{hyperref,listings,float,fancyhdr}
\hypersetup{colorlinks=true,linkcolor=black,citecolor=teal,urlcolor=blue,
  pdftitle={Multimodal Image Retrieval on MIRFLICKR-25K},
  pdfauthor={Ines Besbes and Nesrine Hamrouni}}
\definecolor{codebg}{rgb}{0.97,0.97,0.97}
\definecolor{codeblue}{rgb}{0.13,0.13,0.7}
\definecolor{codegreen}{rgb}{0,0.45,0}
\definecolor{codegray}{rgb}{0.45,0.45,0.45}
\lstset{backgroundcolor=\color{codebg},basicstyle=\ttfamily\footnotesize,
  keywordstyle=\color{codeblue}\bfseries,commentstyle=\color{codegray}\itshape,
  stringstyle=\color{codegreen},breaklines=true,keepspaces=true,numbers=left,
  numberstyle=\tiny\color{codegray},frame=single,rulecolor=\color{gray},
  captionpos=b}
\newcommand{\safefigure}[3]{%
  \begin{figure}[H]\centering
  \IfFileExists{#1}{\includegraphics[width=#2]{#1}}%
    {\fbox{\parbox[c][3cm][c]{0.82\textwidth}{\centering
      Screenshot to add: \texttt{\detokenize{#1}}}}}
  \caption{#3}
  \end{figure}}
\pagestyle{fancy}
\setlength{\headheight}{15pt}
\fancyhf{}
\fancyhead[L]{\small\itshape Multimodal Image Retrieval}
\fancyfoot[C]{\thepage}
\renewcommand{\headrulewidth}{0.4pt}
\begin{document}

\begin{titlepage}
\centering
\vspace*{1.5cm}
{\Large Higher School of Communication of Tunis\par}
\vspace{2cm}
\rule{\linewidth}{0.6mm}\\[0.4cm]
{\huge\bfseries Multimodal Image Retrieval\\on MIRFLICKR-25K\par}
\rule{\linewidth}{0.6mm}\\[2cm]
{\large Project report\par}
\vfill
\begin{minipage}{0.46\textwidth}
\begin{flushleft}\large
\emph{Authors:}\\
\textbf{BESBES In\`es}\\
\textbf{HAMROUNI Nesrine}
\end{flushleft}
\end{minipage}
\hfill
\begin{minipage}{0.46\textwidth}
\begin{flushright}\large
\emph{Faculty adviser:}\\
\textbf{TEBOURBI Riadh}\\[0.5cm]
\emph{Academic year:}\\
\textbf{2026--2027}
\end{flushright}
\end{minipage}
\end{titlepage}

\begin{abstract}
This report presents a multimodal image retrieval system whose primary collection
is MIRFLICKR-25K. It supports visual search from an image, text search over
per-image Flickr tags and topic annotations, and combined image--text search.
The required visual pipeline computes HSV color, LBP texture, HOG shape, and
three MPEG-7 descriptors (CLD, EHD, and SCD). ResNet-18 is an optional
experimental extension in the implementation and is outside the report's
required descriptor scope.
FastAPI serves the application, Elasticsearch indexes text metadata, and
Reciprocal Rank Fusion combines independent text and image lists for OR queries.

The visual index contains 25,000 images and 25,000 signatures for each of the
six required descriptor keys. Text ingestion stores Flickr tags together with
supplied topic annotations. Image-only, text-only, and combined paths have been
checked. These checks confirm that the pipeline runs; they are not a
quantitative evaluation of retrieval quality. Formal precision and recall
measurements on MIRFLICKR annotations remain future work.
\end{abstract}
\noindent\textbf{Keywords:} MIRFLICKR-25K; image retrieval; HSV; LBP; HOG;
MPEG-7; Elasticsearch; fuzzy text search; FastAPI.

\newpage
\tableofcontents
\newpage

\section{Introduction}
\subsection{Motivation and objectives}
Text retrieval can find images whose tags or annotations match a concept.
Content-based image retrieval (CBIR) ranks images by visual similarity to an
example image. The project combines both approaches and supports image-only,
text-only, and combined queries.

This report focuses on MIRFLICKR-25K, which replaces the Objects collection as
the main dataset discussed here. Objects remains supported as a second dataset.
The objectives are to build reusable visual signatures, index available tags
and annotations, implement combined retrieval, and expose the system through a
web interface and REST API.

\section{Dataset}
\subsection{MIRFLICKR-25K collection}
MIRFLICKR-25K is a collection of 25,000 Flickr images introduced for multimedia
retrieval research~\cite{huiskes2008}. In this project, extracted files are
stored under \path{data/mirflickr25k/}. Images are in
\path{data/mirflickr25k/mirflickr/}; topic annotation files are in
\path{data/mirflickr25k/annotations/}; per-image tags are in
\path{data/mirflickr25k/mirflickr/meta/tags/}.

The application uses the full collection, not the older thumbnail directory
\path{data/mirflickr/}. The local scan and visual index both contain 25,000
images.

\subsection{Text metadata}
The ingestion script reads Flickr tags and supplied topic annotation files.
Annotation filenames become topic labels, and image numbers in those files
associate labels with image IDs. Elasticsearch documents store dataset ID,
image ID, filename-derived title, topic words, and tags. EXIF fields are not
currently used as searchable metadata. Tag and annotation coverage varies, so
missing text results do not imply visual irrelevance.

\subsection{Objects collection}
Objects remains available in the application but is not the focus of this
report's counts, screenshots, or evaluation. Results from earlier Objects
experiments must not be presented as MIRFLICKR-25K results.

\section{Data Preparation}
Data preparation has two parts: building visual signatures from the images and
ingesting tags and topic annotations for text search. The MIRFLICKR-25K files
are stored under \path{data/mirflickr25k/}. The collection contains 25,000
images in \path{data/mirflickr25k/mirflickr/}; annotation files and per-image
tag files are stored under \path{data/mirflickr25k/annotations/} and
\path{data/mirflickr25k/mirflickr/meta/tags/}, respectively.

The required visual index uses six descriptor keys: HSV, LBP, HOG, and the
three MPEG-7 descriptors (CLD, EHD, and SCD). Each has 25,000 signatures in the
generated MIRFLICKR index. ResNet-18 remains available in the implementation
as an optional extension but is not included in this report's required
descriptor analysis. Text metadata is ingested separately into Elasticsearch.

\subsection{Dataset Cleansing and Identification}
The image scan reads the 25,000 JPEG files from
\path{data/mirflickr25k/mirflickr/}. Each image is opened and converted to RGB
before visual feature extraction. If an image cannot be read, the indexing
script reports the error and skips that file.

Unlike the Objects collection, MIRFLICKR images are not arranged in
category-specific folders. The image filename, such as \path{im49.jpg},
provides its numeric identifier; it is used to associate the image with its
tag file and any topic annotations. Topic labels come from the supplied
annotation files rather than from a parent directory name.

Image paths are normalized to use forward slashes so identifiers remain
consistent across Windows, Linux, and Docker. The application uses IDs such as
\path{mirflickr/im49.jpg} when matching visual-index entries, Elasticsearch
documents, and image files. During index loading, the dataset prefix is
normalized to align stored feature paths with these application-level IDs.

\section{System Architecture}
Visual preparation scans images, extracts selected descriptors, and writes
vectors to \path{index/mirflickr.pkl}. Text preparation reads Flickr tags and
topic annotations and indexes one Elasticsearch document per image. The visual
index is generated data excluded from Git and must be built from local dataset
files on a new machine.

At query time, FastAPI loads and caches visual signatures on first use. Text
queries go to Elasticsearch. Image queries are encoded with the same descriptor
implementations used during indexing. The API returns ranked results as image
IDs, scores, and URLs.

\begin{figure}[H]
\centering
\fbox{\begin{minipage}{0.88\textwidth}\centering
\textbf{Offline:}\\
MIRFLICKR images $\longrightarrow$ descriptor extraction
$\longrightarrow$ \path{index/mirflickr.pkl}\\[0.25cm]
Flickr tags and topic annotations $\longrightarrow$ Elasticsearch\\[0.4cm]
\textbf{Online:}\\
image and/or text query $\longrightarrow$ visual and/or text retrieval
$\longrightarrow$ AND filtering or OR rank fusion
$\longrightarrow$ API response and web results
\end{minipage}}
\caption{Offline indexing and online retrieval paths.}
\label{fig:architecture}
\end{figure}

\section{Visual Descriptors and Index}
\subsection{Active descriptors}
The report covers six required descriptor keys. HSV, LBP, and HOG use a
$3\times3$ spatial grid, while CLD, EHD, and SCD provide the three MPEG-7
representations. Feature sizes below follow the current code.

\begin{table}[H]
\centering
\small
\begin{tabular}{p{0.25\textwidth}p{0.34\textwidth}rll}
\toprule
\textbf{Index key} & \textbf{Representation} & \textbf{Dimensions}
& \textbf{Distance} & \textbf{Weight}\\
\midrule
\path{hsv_grid3_8_4_4} & HSV histogram, $3\times3$ grid & 1,152 & $\chi^2$ & 0.10\\
\path{lbp_grid3_P8_R1} & Uniform LBP, $P=8$, $R=1$ & 90 & $\chi^2$ & 0.20\\
\path{hog_o9_c16x16_b2x2} & HOG, 9 orientations & 1,764 & $L_2$ & 0.70\\
\path{cld_6_3_3} & MPEG-7 Color Layout (CLD) & 12 & $L_2$ & 0.10\\
\path{ehd_4x4_e5} & MPEG-7 Edge Histogram (EHD) & 85 & $L_1$ & 0.10\\
\path{scd_16_4_4} & MPEG-7 Scalable Color (SCD) & 256 & $L_1$ & 0.10\\
\bottomrule
\end{tabular}
\caption{Active descriptors, dimensions, distances, and configured fusion weights.}
\label{tab:descriptors}
\end{table}

The RGB descriptor class remains in source code, but it is not part of the
active descriptor configuration or MIRFLICKR index. It is not counted as an
available search descriptor here. ResNet-18 is implemented as an optional
experimental feature, but it is excluded from the required descriptor set and
the analysis in this report.

\subsection{Descriptor summary}
\begin{description}
\item[HSV color.] Images are resized to $128\times128$ and divided into nine
cells. Each cell has a normalized $8\times4\times4$ HSV histogram, giving
$9\times128=1{,}152$ components~\cite{swain1991}.
\item[LBP texture.] Grayscale images are resized to $128\times128$. Uniform LBP
histograms with $P=8$ and $R=1$ are computed for each cell. Ten bins per cell
give 90 components~\cite{ojala2002}.
\item[HOG shape.] Grayscale images are resized to $128\times128$. HOG uses nine
orientations, $16\times16$ pixels per cell, and $2\times2$ cells per block,
giving $7\times7\times36=1{,}764$ components~\cite{dalal2005}.
\item[MPEG-7.] CLD represents coarse spatial color with DCT coefficients; EHD
represents edge-type distribution across image regions; SCD represents an HSV
histogram with a Haar transform. The implementations follow the descriptors'
main principles in floating point; they do not claim bit-exact output from
normative reference software~\cite{manjunath2001,iso15938}.
\end{description}

\subsection{Index format}
The pickle file maps descriptor keys to image IDs and feature vectors:
\[
\mathrm{Index}=\{\mathrm{descriptor}:
\{\mathrm{imageID}\mapsto\mathbf{v}\in\mathbb{R}^{D}\}\}.
\]
The report-scope index has 25,000 vectors for each of six required keys, or
150,000 signatures. The local pickle may also contain the optional ResNet
signatures, which are excluded from the report's descriptor count. The API
stacks vectors into aligned NumPy matrices. The file is
\path{index/mirflickr.pkl}; it is locally generated and not committed to Git.

\section{Retrieval and Fusion}
\subsection{Image-only search}
For descriptor $d$, the system computes query distances to all candidates.
Distances are min--max normalized over the candidate set:
\[
\widehat{\delta}_d(q,x)=
\frac{\delta_d(q,x)-\min_x\delta_d(q,x)}
{\max_x\delta_d(q,x)-\min_x\delta_d(q,x)}.
\]
The weighted score is
\[
S(q,x)=\frac{\sum_{d\in D}w_d\,\widehat{\delta}_d(q,x)}
{\sum_{d\in D}w_d},
\]
and images are returned in ascending score order. A zero-range vector is
protected against division by zero.

The configured weights are not learned from a MIRFLICKR relevance study.
MIRFLICKR calibration is needed before treating these settings as optimal.
Current ranking does not use the category-vote reranking described in the
earlier Objects-focused report. The report's descriptor analysis uses the six
required descriptors; the optional ResNet extension is excluded.

\subsection{Text-only search}
Elasticsearch searches \path{category_words}, \path{tags}, and \path{title}.
Topic words receive a field boost of 3, tags a boost of 2, and title its default
boost. The analyzer lowercases text, folds accents, and applies the Porter
stemmer. The query uses \path{fuzziness=AUTO:4,8} for typo tolerance. Results
are restricted to the selected dataset.

\subsection{Combined search}
\begin{itemize}
\item \textbf{AND:} Elasticsearch obtains up to 200 text candidates. Visual
search ranks those candidates with the selected descriptors.
\item \textbf{OR:} Text and visual search run independently; their lists are
combined with Reciprocal Rank Fusion (RRF)~\cite{cormack2009}:
\[
\mathrm{RRF}(x)=
\sum_{m\in\{\mathrm{text},\mathrm{visual}\}}\frac{1}{60+r_m(x)},
\]
where $r_m(x)$ is the one-based rank of image $x$ in modality $m$.
\end{itemize}
RRF combines ranks rather than raw Elasticsearch scores and visual distances.
It is used for OR fusion; AND uses text results as candidates for visual ranking.

\section{Text Indexing, API, and Web Client}
Each MIRFLICKR Elasticsearch document contains image ID, dataset ID, URL,
filename-derived title, topic labels, and Flickr tags. Tags and annotations are
complementary, but noisy or incomplete; they are not exhaustive captions. The
default Elasticsearch index name is \path{cbir_images}.

FastAPI serves search, image, dataset-status, and catalog routes. The main
endpoints are \path{GET /health}, \path{GET /datasets}, \path{GET /catalog},
\path{POST /search}, \path{POST /search/upload}, and \path{GET /images/...}.
The web interface lets users select a dataset, choose visual descriptor
families, set result count, enter keywords, choose or upload an image, and
start another query from a result.

\safefigure{screenshots/mirflickr-interface.png}{0.88\textwidth}
{MIRFLICKR-25K search interface. Add a current screenshot at this path.}
\safefigure{screenshots/mirflickr-results.png}{0.88\textwidth}
{Example MIRFLICKR-25K results. Add a current results screenshot at this path.}

\section{Reproducibility and Deployment}
The following commands prepare the text and visual indexes and start the API.
These commands build the six descriptors required for the report.
\begin{lstlisting}[language=bash,caption={Local MIRFLICKR setup}]
python -m venv .venv
.\.venv\Scripts\Activate.ps1
pip install -r requirements.txt

docker compose up -d elasticsearch
python scripts/ingest_es.py --dataset mirflickr

python scripts/build_index.py --dataset mirflickr \
  --descriptors hsv lbp hog cld ehd scd

python -m uvicorn api.main:app --reload --port 8000
\end{lstlisting}
Open \path{http://localhost:8000/web/}; API documentation is at
\path{http://localhost:8000/docs}. If port 8000 is unavailable, use another port
such as \path{--port 8766} and update the browser URL accordingly.

The repository includes Docker and Compose configuration for FastAPI and
Elasticsearch. Data and generated index files are mounted into the API
container. ResNet-18 remains an optional implementation extension and needs
PyTorch and Torchvision, which the current Dockerfile does not install. The
generated index is ignored by Git; build it on each new installation and
restart the API so its in-memory cache reloads the index.

\section{Verification, Limitations, and Future Work}
\subsection{Implementation checks}
The local MIRFLICKR index reports 25,000 images. Each of the six required
descriptor keys has 25,000 signatures, for 150,000 signatures in the
report-scope descriptor set. The local file may also contain the optional
ResNet signatures. The following checks describe index coverage and software
functionality, not retrieval relevance:

\begin{table}[H]
\centering
\small
\begin{tabular}{p{0.31\textwidth}p{0.27\textwidth}p{0.32\textwidth}}
\toprule
\textbf{Check} & \textbf{Status} & \textbf{Meaning}\\
\midrule
Images in visual collection & 25,000 & Collection coverage\\
Signatures per required descriptor & 25,000 & Six descriptors indexed\\
Total required signatures & 150,000 & Excludes optional ResNet\\
Image-only, text-only, AND, OR & Smoke-checked & Confirms query paths run\\
Precision@10, recall, mAP & Not measured & No MIRFLICKR quality benchmark yet\\
\bottomrule
\end{tabular}
\caption{Implementation checks and retrieval evaluation status.}
\label{tab:evaluation-status}
\end{table}

The Objects precision figures from the earlier report are not MIRFLICKR
measurements and are not carried over here.

\subsection{Limitations}
\begin{itemize}
\item No MIRFLICKR precision@k, recall, mean average precision, or NDCG has yet
been measured. Objects measurements must not be reused as MIRFLICKR results.
\item Flickr tags are user supplied and annotation coverage varies. EXIF is
not indexed by the current text pipeline.
\item HSV, LBP, HOG, and MPEG-7 can match background, color, or texture instead
of the main subject.
\item Descriptor weights have not been calibrated with MIRFLICKR relevance
judgments. Min--max normalized distances depend on the candidate set and are
not probabilities.
\end{itemize}

\subsection{Future work}
Define a reproducible query set from MIRFLICKR topic annotations and compare
individual descriptors with fusion. Report precision@k, recall@k, and mean
average precision for image-only search, and evaluate text-only and AND/OR
search separately. Calibrate weights on a development split and report results
on a held-out split. Further extensions include stronger image embeddings,
evaluation of the optional ResNet-18 implementation~\cite{he2016}, improved tag
normalization, optional EXIF indexing, and approximate-nearest-neighbor search
for larger collections.

\section{Conclusion}
MIRFLICKR-25K is the primary collection in this report. Its visual index
contains HSV, LBP, HOG, and MPEG-7 CLD/EHD/SCD signatures for all 25,000 images.
The implementation also retains ResNet-18 as an optional extension outside
the required scope. The application combines visual search with Flickr tags and
topic annotations, supporting image-only, text-only, and combined retrieval.
Text queries use Elasticsearch fuzzy matching, and OR queries use RRF.

The implementation is operational, but MIRFLICKR quality has not been
established with a formal relevance benchmark. Evaluation against supplied
topic annotations is needed before making quantitative accuracy claims.

\begin{thebibliography}{99}
\bibitem{huiskes2008}
M.~J. Huiskes and M.~S. Lew,
\emph{The MIR Flickr Retrieval Evaluation},
in \emph{Proceedings of the 1st ACM International Conference on Multimedia
Information Retrieval}, 2008, pp.~39--43.
doi: \href{https://doi.org/10.1145/1460096.1460104}
{10.1145/1460096.1460104}.

\bibitem{swain1991}
M.~J. Swain and D.~H. Ballard, \emph{Color indexing},
\emph{International Journal of Computer Vision}, vol.~7, no.~1,
pp.~11--32, 1991.

\bibitem{ojala2002}
T.~Ojala, M.~Pietik{\"a}inen, and T.~M{\"a}enp{\"a}{\"a},
\emph{Multiresolution gray-scale and rotation invariant texture classification
with local binary patterns},
\emph{IEEE Transactions on Pattern Analysis and Machine Intelligence},
vol.~24, no.~7, pp.~971--987, 2002.

\bibitem{dalal2005}
N.~Dalal and B.~Triggs,
\emph{Histograms of oriented gradients for human detection},
in \emph{IEEE Conference on Computer Vision and Pattern Recognition}, 2005,
pp.~886--893.

\bibitem{manjunath2001}
B.~S. Manjunath, J.-R. Ohm, V.~V. Vasudevan, and A.~Yamada,
\emph{Color and texture descriptors},
\emph{IEEE Transactions on Circuits and Systems for Video Technology},
vol.~11, no.~6, pp.~703--715, 2001.

\bibitem{iso15938}
ISO/IEC 15938-3, \emph{Information Technology---Multimedia Content Description
Interface---Part 3: Visual}, 2002.

\bibitem{he2016}
K.~He, X.~Zhang, S.~Ren, and J.~Sun,
\emph{Deep residual learning for image recognition},
in \emph{IEEE Conference on Computer Vision and Pattern Recognition}, 2016,
pp.~770--778.

\bibitem{cormack2009}
G.~V. Cormack, C.~L. Clarke, and B.~B{\"u}ttcher,
\emph{Reciprocal rank fusion outperforms Condorcet and individual rank learning
methods},
in \emph{Proceedings of the 32nd ACM SIGIR Conference}, 2009,
pp.~758--759.
\end{thebibliography}
\end{document}
